\documentclass[11pt]{article}

\usepackage{amsmath}
\usepackage{amssymb}
\usepackage{amsthm}
\usepackage[round]{natbib}
\usepackage{graphicx}
\usepackage{hyperref}
\usepackage{booktabs}
\usepackage{array}

\newtheorem{theorem}{Theorem}
\newtheorem{lemma}[theorem]{Lemma}
\newtheorem{corollary}[theorem]{Corollary}

\hypersetup{
  colorlinks=true,
  linkcolor=blue,
  citecolor=blue,
  urlcolor=blue,
}

\title{A Four-Prong Obstruction to $(3{+}1)$-Dimensional Emergence
       from $\mathrm{Cl}(6)$}
\author{Todd M.~Firestone\\
        Independent Researcher}
\date{June 2026}

\begin{document}
\maketitle

%======================================================================
% ABSTRACT
%======================================================================
\begin{abstract}
We prove that the $\mathrm{Cl}(6)$ bivector algebra $\mathfrak{so}(6)$,
viewed as the internal symmetry algebra of a $\mathrm{Cl}(6)$-based
emergence programme, admits no clean assignment of spatial rotation
generators compatible with the simultaneous preservation of
$\mathfrak{su}(3)$ color.  The obstruction is dimension-theoretic and
basis-independent: any designation of two internal directions as spatial
leaves a four-direction complement whose bivector algebra is
$\mathfrak{so}(4)$, of dimension six, which is strictly less than the
eight required to house $\mathfrak{su}(3)$.  Combined with three further
algebraic obstructions---the one-dimensional centralizer of
$\mathfrak{su}(3)$ in $\mathfrak{so}(6)$, the one-dimensional nature of
the chain axis that supplies the $1{+}1$ sector, and the parallel
double-booking dilemma on the quaternion factor
$\mathbb{H}$---the result constitutes a four-prong structural no-go for
$(3{+}1)$ emergence under the full Stance-1 hypothesis (no imported
spacetime or gauge structure).  All four prongs are computationally
verified at machine precision.  The no-go sharpens and completes earlier
partial results.  Under Stance~1 the only surviving escape is whether the
emergent infrared theory voids the hypotheses of the Coleman--Mandula
theorem by being genuinely gapless with no asymptotic particle states;
this remains open.  Under Stance~2, the centralizer theorem provides an
algebraic justification for every import rather than a default assumption.
\end{abstract}

%======================================================================
% SECTION 1: INTRODUCTION
%======================================================================
\section{Introduction}

The TMFT emergence programme asks whether a $(3{+}1)$-dimensional
spacetime containing the Standard Model gauge group can be derived,
without import, from the Clifford algebra $\mathrm{Cl}(6)$ tensored with
$\mathrm{Cl}(1,3)$.  The C-series lattice construction has successfully
produced a clean $1{+}1$ Lorentzian sector: a massless chiral fermion
along an external chain axis~$n$, with the Tomita--Takesaki modular
structure supplying the $1{+}1$ Lorentzian
signature~\citep{C4consolidation}.  The $1{+}1$ result is established.
The question addressed here is whether the same construction, or any
algebraic extension of it within the same $\mathrm{Cl}(6)$ structure, can
produce two additional spatial directions and recover $(3{+}1)$.
% [ESTABLISHED]

The $\mathrm{Cl}(6)$ bivector algebra is $\mathfrak{so}(6)$, of
dimension~15.  Its maximal $\mathfrak{su}(3)$-color subalgebra occupies
eight of those fifteen dimensions.  A prior
result~\citep{centralizercheck} established that the centralizer of
$\mathfrak{su}(3)_{\mathrm{color}}$ in $\mathfrak{so}(6)$ is exactly
one-dimensional, spanned by $H = B_{12} + B_{34} + B_{56}$, which is
proportional to the baryon-minus-lepton charge $Q_{B\text{-}L}$.  That
single result has four distinct structural consequences, each closing a
different candidate route to $(3{+}1)$.  Together they form a complete
no-go theorem.

The paper is organized as follows.  Section~\ref{sec:centralizer} states
and verifies the centralizer theorem.  Section~\ref{sec:consequences}
derives the four consequences.  Section~\ref{sec:partition} states the
partition-double-booking dilemma and its dimension-theoretic proof.
Section~\ref{sec:quaternion} closes the quaternion spatial route.
Section~\ref{sec:nogo} assembles the four-prong no-go.
Section~\ref{sec:survives} identifies the two surviving paths: the beta
question (whether the emergent infrared voids Coleman--Mandula's
hypotheses by being genuinely gapless with no asymptotic particle states)
and Stance~2 (imported spacetime and gauge structure, the fallback if
beta fails, with each import retrospectively justified by the centralizer
theorem rather than assumed by default).
Section~\ref{sec:discussion} discusses the structure of the no-go and why
all four prongs are necessary.

\paragraph{Conventions.}
$\mathrm{Cl}(6)$ is generated by $e_1, \ldots, e_6$ with
$\{e_i, e_j\} = 2\,\delta_{ij}$, all $e_i$ Hermitian.  Bivectors
$B_{ij} = \tfrac{1}{2}[e_i, e_j]$ span $\mathfrak{so}(6)$.  Throughout,
the candidate symmetry algebra is the Lie algebra of bivectors
$\mathfrak{so}(6)$, which generates the spin group $\mathrm{Spin}(6)$ by
exponentiation.  This is the standard choice in geometric-algebra
physics: gauge and spacetime symmetries act through the bivector algebra.
Higher-grade elements of $\mathrm{Cl}(6)$ (4-vectors, 6-vectors, the
volume element), when exponentiated, generate extended groups (conformal
or higher-graded transformations) outside the
$\mathrm{Spin}(p,q)$ gauge and spacetime framework of the present
construction, and are excluded on those grounds.  The $\mathfrak{su}(3)$
subalgebra is derived as the stabilizer of the $H = -3/2$ eigenstate of
the number operator.  The stabilizer construction of $\mathfrak{su}(3)$
inside $\mathfrak{so}(6) \cong \mathfrak{su}(4)$ goes back
to~\citet{GunaydinGursey1973} and was elaborated
by~\citet{Dixon1994}; the formulation used here
follows~\citet{Furey2016}.
% [INHERITED]
The quaternion factor $\mathbb{H}$ acts on a four-dimensional real space
($\mathbb{H}$ as real vector space) via left and right multiplication;
$\mathfrak{su}(2)_L = \operatorname{span}\{L_i/2, L_j/2, L_k/2\}$ and
$\mathfrak{su}(2)_R = \operatorname{span}\{-R_i/2, -R_j/2, -R_k/2\}$.

%======================================================================
% SECTION 2: THE CENTRALIZER THEOREM
%======================================================================
\section{The Centralizer Theorem}
\label{sec:centralizer}

\begin{theorem}[Centralizer]
\label{thm:centralizer}
The centralizer of\/ $\mathfrak{su}(3)_{\mathrm{color}}$ inside the\/
$\mathrm{Cl}(6)$ bivector Lie algebra $\mathfrak{so}(6)$ is exactly
one-dimensional:
\begin{equation}
\label{eq:centralizer}
C_{\mathfrak{so}(6)}\!\bigl(\mathfrak{su}(3)\bigr)
  = \operatorname{span}\{H\},
\qquad
H = B_{12} + B_{34} + B_{56} = \tfrac{3}{2}\,Q_{B\text{-}L}.
\end{equation}
No other bivector commutes with all eight color generators.
\end{theorem}
% [ESTABLISHED -- Schur argument below; corroborated in rep by cl6_su3_centralizer_check.py]

\begin{proof}
The result follows from Schur's lemma, basis-independently.  The
$\mathrm{Cl}(6)$ bivector algebra is
$\mathfrak{so}(6) \cong \mathfrak{su}(4)$, acting on $\mathbb{C}^4$.
The $\mathfrak{su}(3)_{\mathrm{color}}$ subalgebra is the Furey embedding
$\mathfrak{su}(3) \subset \mathfrak{su}(4)$, under which $\mathbb{C}^4$
decomposes as the reducible $\mathbf{3} \oplus \mathbf{1}$: the
fundamental $\mathbf{3}$ of $\mathfrak{su}(3)$ plus a singlet.  By
Schur's lemma, an element of $\mathrm{End}(\mathbb{C}^4)$ commuting with
all of $\mathfrak{su}(3)$ acts as a scalar on each irreducible block, so
the commutant of $\mathfrak{su}(3)$ in $\mathrm{End}(\mathbb{C}^4)$ is
$\mathbb{C} \oplus \mathbb{C}$, two complex-dimensional, spanned by the
two block projectors.  Restricting to $\mathfrak{su}(4)$ (traceless
Hermitian, in the physics convention where $\mathfrak{su}(4)$ generators
are Hermitian) reduces the four real dimensions of the
$\mathbb{C} \oplus \mathbb{C}$ commutant to one: Hermiticity forces
$\alpha$ and $\beta$ real (removing two), and tracelessness imposes
$3\alpha + \beta = 0$ (removing one), leaving one real dimension.  The
surviving generator is proportional to $\operatorname{diag}(1,1,1,-3)$,
which is exactly the hypercharge\,/\,$Q_{B\text{-}L}$ assignment.  The
dimension of the centralizer is independent of basis and of
representation; the identification with the specific bivector combination
$H = B_{12} + B_{34} + B_{56}$ fixes the bivector basis to the canonical
mode-pair decomposition $(e_1 e_2)$, $(e_3 e_4)$, $(e_5 e_6)$.
\end{proof}

The script \texttt{cl6\_su3\_centralizer\_check.py} corroborates this in
the explicit 8-dimensional spinor representation: the centralizer kernel
is one-dimensional with alignment
$|\langle v, H\rangle|/\|v\|\,\|H\| = 1.0000000000$ at machine
precision, and the $\mathfrak{su}(3)$ octet closes at residual
$6.2 \times 10^{-16}$.

\begin{lemma}[Decomposition]
\label{lem:decomposition}
Every bivector decomposes across three orthogonal\/ $\mathfrak{su}(3)$
sectors---adjoint-$8$, centralizer-$1$, and coset-$(\mathbf{3} +
\bar{\mathbf{3}})$---as follows:
\begin{itemize}
\item Each within-pair Cartan element $B_{12}$, $B_{34}$, $B_{56}$ has
  exactly $2/3$ adjoint content and $1/3$ centralizer content, zero coset.
\item Each cross-plane bivector has exactly $1/2$ adjoint content and
  $1/2$ coset content, zero centralizer.
\item No bivector lies purely in any single sector.
\end{itemize}
\end{lemma}
% [ESTABLISHED -- projection argument below; corroborated in rep by cl6_spatial_direction_carving.py, Part A]

The fractions in Lemma~\ref{lem:decomposition} are not numerical
accidents.  The three-dimensional Cartan space
$\operatorname{span}\{B_{12}, B_{34}, B_{56}\}$ decomposes as
$\operatorname{span}\{h_1, h_2\}$ (two-dimensional $\mathfrak{su}(3)$
Cartan) $\oplus$ $\operatorname{span}\{H\}$ (one-dimensional
centralizer); each Cartan bivector therefore carries $2/3$ adjoint and
$1/3$ centralizer by a simple projection count.  The cross-plane result
follows from the absence of $H$-content in off-diagonal bivectors, which
is forced by the coset structure.

\begin{corollary}
\label{cor:uniqueU1}
Within $\mathfrak{so}(6)$, exactly one independent $\mathrm{U}(1)$
commutes with $\mathfrak{su}(3)_{\mathrm{color}}$---the one generated by
$H = \tfrac{3}{2}\,Q_{B\text{-}L}$.  Any structure requiring more than
one commuting generator, or requiring a commuting generator free of
$B$--$L$ charge, cannot be housed in the $\mathrm{Cl}(6)$ bivector
algebra.
\end{corollary}
% [ESTABLISHED -- immediate from Theorem 1]

%======================================================================
% SECTION 3: FOUR CONSEQUENCES
%======================================================================
\section{Four Consequences of One Theorem}
\label{sec:consequences}

The single fact that
$\dim C_{\mathfrak{so}(6)}\!\bigl(\mathfrak{su}(3)\bigr) = 1$ closes
four independent candidate assignments.  We take each in turn.

\subsection{Hypercharge---the seat that fits}

$\mathrm{U}(1)_Y$ hypercharge requires one generator commuting with
$\mathfrak{su}(3)_{\mathrm{color}}$.
Theorem~\ref{thm:centralizer} provides exactly one: $H$.  Identifying
$Y \propto H$ (Paper~2: $Q_{B\text{-}L} = (2/3)\,H$) places hypercharge
in the unique color-neutral bivector.  The centralizer being
one-dimensional explains why there is precisely one independent commuting
$\mathrm{U}(1)$---not zero, not two.  Hypercharge takes that seat
completely.
% [ESTABLISHED]

\subsection{Weak isospin $\mathrm{SU}(2)_L$---locked out of $\mathrm{Cl}(6)$ bivectors}

$\mathrm{SU}(2)_L$ requires three generators commuting with
$\mathfrak{su}(3)_{\mathrm{color}}$.  The centralizer has dimension one.
Three independent generators cannot fit.
% [ESTABLISHED]

This provides the algebraic resolution of the B-2 finding from the
Paper~3 adversarial referee review (June~2026): the corrected
$T_3 = -B_{34}$ fails to commute with the color $\mathfrak{su}(3)$
because $B_{34}$ is $2/3$ color-adjoint
(Lemma~\ref{lem:decomposition}).  The B-2 finding is the centralizer
theorem asserting itself on the electroweak sector.  No choice of
$T_1, T_2, T_3$ from within the $\mathrm{Cl}(6)$ bivector algebra can
commute with all of $\mathfrak{su}(3)_{\mathrm{color}}$.
$\mathrm{SU}(2)_L$ must be accommodated on the separate $\mathbb{H}$
factor of
$\mathbb{R} \otimes \mathbb{C} \otimes \mathbb{H} \otimes \mathbb{O}$---the
standard Furey construction~\citep{Furey2016}.

\subsection{Spatial $\mathfrak{so}(3)$---locked out (two mechanisms)}

Spatial rotation $\mathfrak{so}(3)$ requires three generators.  The
centralizer route needs
$\dim \geq 3$ in
$C_{\mathfrak{so}(6)}\!\bigl(\mathfrak{su}(3)\bigr)$;
Theorem~\ref{thm:centralizer} gives $\dim = 1$.  This closes the
double-booking route.  The partition route is closed separately and is the
subject of Section~\ref{sec:partition}.

\subsection{Color-neutral photon helicity---locked out ($B$--$L$ reductio)}

A color-neutral helicity generator would have to be proportional to $H$
(Corollary~\ref{cor:uniqueU1}).  But
$H = \tfrac{3}{2}\,Q_{B\text{-}L}$ carries baryon-minus-lepton charge.
Photons carry $B$--$L = 0$.  The contradiction is immediate: transverse
polarization modes cannot be color-neutral within the $\mathrm{Cl}(6)$
bivector algebra.  This is the algebraic half of the two-walls identity
established in~\citep{TwoWalls}.
% [ESTABLISHED -- B-L reductio; alignment corroborated in rep, cl6_su3_centralizer_check.py: 1.0000000000]

\subsection{Summary}

\begin{table}[ht]
\centering
\begin{tabular}{@{}l c c l@{}}
\toprule
Sector & Generators needed &
  Fits in centralizer ($\dim 1$)? & Verdict \\
\midrule
$\mathrm{U}(1)_Y$ hypercharge & 1 &
  Yes---$Y \propto H$ & FITS \\
$\mathrm{SU}(2)_L$ weak isospin & 3 &
  No---$3 > 1$ & LOCKED OUT \\
Spatial $\mathfrak{so}(3)$ & 3 &
  No---$3 > 1$ (double-booking) & LOCKED OUT \\
Color-neutral helicity & 1, $B$--$L$-free &
  No---$H$ carries $B$--$L$ & LOCKED OUT \\
\bottomrule
\end{tabular}
\caption{Centralizer assignments.  One seat.  Four candidates.  Only
  hypercharge fits.}
\label{tab:centralizer-summary}
\end{table}

%======================================================================
% SECTION 4: THE PARTITION-DOUBLE-BOOKING DILEMMA
%======================================================================
\section{The Partition-Double-Booking Dilemma}
\label{sec:partition}

The C-series construction produces a $1{+}1$ Lorentzian sector along an
external chain axis~$n$~\citep{C4consolidation}.  The chain axis is
external---it does not come from the $\mathrm{Cl}(6)$ internal
directions.  For $(3{+}1)$, two additional spatial directions must be
produced.  The only candidate source within the current construction is
the six-dimensional space spanned by $e_1, \ldots, e_6$.  We now show
that no assignment of two of these directions as spatial is consistent
with the survival of $\mathfrak{su}(3)_{\mathrm{color}}$.

\begin{theorem}[Partition No-Go]
\label{thm:partition}
Let $V_2$ be any two-dimensional subspace of\/ $\mathbb{R}^6$.  Its
orthogonal complement $V_4$ has bivector algebra $\mathfrak{so}(V_4)$,
with
\begin{equation}
\label{eq:partition-dim}
\dim\mathfrak{so}(V_4) = \binom{4}{2} = 6
  < 8 = \dim\mathfrak{su}(3).
\end{equation}
Therefore no eight-dimensional Lie algebra, $\mathfrak{su}(3)$ in
particular, is a subalgebra of\/ $\mathfrak{so}(V_4)$.  Under any
partition of the six\/ $\mathrm{Cl}(6)$ directions into two spatial and
four internal, $\mathfrak{su}(3)_{\mathrm{color}}$ cannot survive as a
subalgebra of the residual internal bivector algebra.
\end{theorem}
% [ESTABLISHED -- basis-independent combinatorial proof; corroborated in rep for all cases, cl6_spatial_direction_carving.py]

\begin{proof}
The bivector count is purely combinatorial.  The dimension of
$\mathfrak{so}(V_4)$ is the number of 2-element subsets of a 4-element
basis: $\binom{4}{2} = 6$.  Since $\mathfrak{su}(3)$ has dimension~8
and $8 > 6$, no Lie algebra isomorphic to $\mathfrak{su}(3)$ can be a
subalgebra of the six-dimensional $\mathfrak{so}(V_4)$.  This is
independent of the choice of $V_2$.
\end{proof}

\paragraph{Numerical concretization.}
Theorem~\ref{thm:partition} is the basis-independent statement; the
script \texttt{cl6\_spatial\_direction\_carving.py} concretizes it in a
fixed basis by recording, for each partition, how many of the eight
original $\mathfrak{su}(3)$ generators remain intact (free of
spatial-sector components) in the residual.  This is a stronger per-case
fact than the theorem requires---the theorem only needs that no
$\mathfrak{su}(3)$ subalgebra exists in $\mathfrak{so}(V_4)$---but the
per-case leakage count corroborates it concretely.  The script tests all
inequivalent assignment types (two cases up to $S_3$ permutation symmetry
of the mode pairs, plus all eight cross-pair variants exhaustively):

\begin{itemize}
\item \textbf{Case~A (within-pair):} spatial directions $e_5, e_6$.
  Spatial rotation generator $B_{56}$.  Six purely internal bivectors
  span $\mathfrak{so}(4)$ at closure residual $3.14 \times 10^{-16}$;
  $\mathfrak{su}(3)$ survival $= 1/8$ generators intact.  $H$ disrupted
  at $\|H - \mathrm{proj}\|/\|H\| = 0.577$.
  \textbf{$\mathfrak{su}(3)$ broken.}

\item \textbf{Case~B (cross-pair):} spatial directions $e_3, e_5$.
  Spatial rotation generator $B_{35}$.  Six purely internal bivectors
  span $\mathfrak{so}(4)$ at closure residual $3.14 \times 10^{-16}$;
  $\mathfrak{su}(3)$ survival $= 0/8$ generators intact.  $H$ disrupted
  at $\|H - \mathrm{proj}\|/\|H\| = 0.816$.
  \textbf{$\mathfrak{su}(3)$ broken.}

\item \textbf{All eight cross-pair variants:} every assignment yields six
  internal bivectors, dimension six, $0/8$ $\mathfrak{su}(3)$ generators
  surviving.  The result is uniform across all choices.
\end{itemize}

\paragraph{The dilemma.}
The double-booking obstruction (Sections~\ref{sec:centralizer}--\ref{sec:consequences}, consequence~(iii)) and the partition no-go
(Theorem~\ref{thm:partition}) together form a dilemma with no middle
ground:

\begin{itemize}
\item \textbf{If you partition} (designate two directions spatial, keep
  four internal): $\mathfrak{su}(3)$ is destroyed in the residual sector
  by Theorem~\ref{thm:partition}.
\item \textbf{If you do not partition} (all six directions serve dual
  spatial and color roles---double-booking): spatial rotations
  $\mathfrak{so}(3)$ fail to commute with color $\mathfrak{su}(3)$, since
  $C_{\mathfrak{so}(6)}\!\bigl(\mathfrak{su}(3)\bigr) =
  \operatorname{span}\{H\}$ is one-dimensional and cannot house the three
  generators $\mathfrak{so}(3)$ requires
  (Theorem~\ref{thm:centralizer}, consequence~(iii)).
\end{itemize}

The six internal directions are the only reservoir.  Either partition them
or do not.  Each horn is fatal by a different mechanism.
% [ESTABLISHED reasoning]

%======================================================================
% SECTION 5: THE QUATERNION SPATIAL ROUTE
%======================================================================
\section{The Quaternion Spatial Route---Closed}
\label{sec:quaternion}

The dilemma of Section~\ref{sec:partition} operates entirely within the
$\mathrm{Cl}(6)$ bivector algebra.  A natural candidate escape is to seek
spatial rotation generators from the quaternion factor $\mathbb{H}$ of
$\mathbb{R} \otimes \mathbb{C} \otimes \mathbb{H} \otimes \mathbb{O}$,
particularly since $\mathrm{SU}(2)_L$ is already assigned there
(Section~\ref{sec:consequences}).  We now show that $\mathbb{H}$ provides
no clean spatial route either.

\begin{theorem}[$\mathbb{H}$ Dilemma]
\label{thm:Hdilemma}
The quaternion algebra $\mathbb{H}$, viewed as a four-dimensional real
vector space with left and right multiplication algebras
\begin{align}
\mathfrak{su}(2)_L &= \operatorname{span}\{L_i/2,\, L_j/2,\, L_k/2\},
\\
\mathfrak{su}(2)_R &= \operatorname{span}\{-R_i/2,\, -R_j/2,\,
  -R_k/2\},
\end{align}
satisfies:
\begin{align}
C_{\operatorname{span}\{L_a, R_b, I\}}
  \!\bigl(\mathfrak{su}(2)_L\bigr)
  &= \mathfrak{su}(2)_R \oplus \mathbb{R}
  \qquad (\dim 4),
\label{eq:cent-suL}
\\
C_{\operatorname{span}\{L_a, R_b, I\}}
  \!\bigl(\mathfrak{su}(2)_L \oplus \mathfrak{su}(2)_R\bigr)
  &= \operatorname{span}\{I\}
  \qquad (\dim 1).
\label{eq:cent-full}
\end{align}
The scalar identity $I$ is not a rotation generator.  Therefore no
non-trivial generator in $\mathbb{H}$ commutes with both internal\/
$\mathrm{SU}(2)$ factors.
\end{theorem}
% [ESTABLISHED -- forced by associativity of H; corroborated in rep, cl6_H_spatial_route_check.py, all 6 checks 0.00e+00]

\begin{proof}
Direct computation (script
\texttt{cl6\_H\_spatial\_route\_check.py}):

\begin{itemize}
\item $\mathfrak{su}(2)_L$ closes:
  $[L_a, L_b] = 2\,\varepsilon_{abc}\,L_c$,
  residual $0.00 \times 10^{+00}$.
  \textsc{[Check~A]}

\item $\mathfrak{su}(2)_R$ closes:
  $[R_a, R_b] = -2\,\varepsilon_{abc}\,R_c$,
  residual $0.00 \times 10^{+00}$;
  $S_a = -R_a/2$ satisfies
  $[S_a, S_b] = \varepsilon_{abc}\,S_c$ (genuine $\mathfrak{su}(2)$).
  \textsc{[Check~B]}

\item Cross-commutativity:
  $[L_a, R_b] = 0$ for all nine pairs, $0.00 \times 10^{+00}$.
  \textsc{[Check~C]}

\item Centralizer of $\mathfrak{su}(2)_L$: basis
  $\{I, R_i, R_j, R_k\}$, $\dim 4$, zero $L_a$ content.
  \textsc{[Check~D]}

\item $\mathfrak{su}(2)_R$ is nonabelian:
  $\max\|[R_a, R_b]\| = 4.0$ (not a $\mathrm{U}(1)$).
  \textsc{[Check~E]}

\item Centralizer of $\mathfrak{su}(2)_L \oplus \mathfrak{su}(2)_R$:
  basis $\{I\}$, $\dim 1$, identity fraction $1.000000$.
  \textsc{[Check~F]}
\end{itemize}
\end{proof}

\paragraph{The $\mathbb{H}$ dilemma.}
The quaternion factor $\mathbb{H}$ is the internal factor of the
$\mathbb{R} \otimes \mathbb{C} \otimes \mathbb{H} \otimes \mathbb{O}$
decomposition by construction; every element of its rotation algebra
$\mathfrak{so}(4) = \mathfrak{su}(2)_L \oplus \mathfrak{su}(2)_R$ is
therefore internal, including $\mathfrak{su}(2)_L$ (weak isospin),
$\mathfrak{su}(2)_R$, the diagonal
$\mathfrak{su}(2)_{\mathrm{diag}} =
\operatorname{span}\{(L_a - R_a)/2\}$, and any other three-dimensional
subalgebra.  A spatial rotation $\mathfrak{so}(3)$ would have to commute
with the entire internal $\mathbb{H}$-structure to avoid double-booking.
The load-bearing fact is the centralizer:
$C\!\bigl(\mathfrak{su}(2)_L \oplus \mathfrak{su}(2)_R\bigr) =
\operatorname{span}\{I\}$ [Check~F].  Any element of the $\mathbb{H}$
rotation algebra that commutes with the full internal structure must be
proportional to the scalar identity $I$, and $I$ is not a rotation
generator ($\exp(\theta\,I)$ is a dilation of $\mathbb{R}^4$, not a
rotation; $I$ is symmetric, not antisymmetric, hence not in
$\mathfrak{so}(4)$).  No three-dimensional spatial rotation algebra can
be sourced from~$\mathbb{H}$.

This subsumes the individual subalgebra cases.  The candidate
$\mathfrak{su}(2)_{\mathrm{diag}}$, for instance, does not commute with
$\mathfrak{su}(2)_L$:
$[L_a - R_a,\, L_b] = [L_a, L_b] = 2\,\varepsilon_{abc}\,L_c$ is
nonzero for $a \neq b$ (using cross-commutativity $[R_a, L_b] = 0$,
Check~C).  The same computation rules out any other would-be spatial
subalgebra by the centralizer result directly, with no need to enumerate
cases.

The obstruction is algebraic and complete at this level.  The
Coleman--Mandula confrontation enters only conditionally, as in
prong~(ii): if the emergence programme reaches a 4D theory meeting
Coleman--Mandula's hypotheses (Section~\ref{sec:survives}), mixing
Poincar\'{e} spatial rotations with any internal symmetry on $\mathbb{H}$
is forbidden there as well, independently of the tensor factor on which
the mixing occurs.  The algebraic obstruction does not need
Coleman--Mandula; Coleman--Mandula is the backup that closes the same door
at the level of the target theory~\citep{ColemanMandula1967}.
% [ESTABLISHED reasoning]

%======================================================================
% SECTION 6: THE FOUR-PRONG NO-GO
%======================================================================
\section{The Four-Prong No-Go}
\label{sec:nogo}

We assemble the results of the preceding sections.

\begin{theorem}[Four-Prong No-Go]
\label{thm:nogo}
Under Stance~1 (no imported spacetime or gauge structure), the\/
$\mathbb{R} \otimes \mathbb{C} \otimes \mathbb{H} \otimes \mathbb{O}$
algebraic construction admits no clean $(3{+}1)$-dimensional extension of
the established $1{+}1$ sector.  Four independent obstructions hold
simultaneously:
\end{theorem}
% [ESTABLISHED reasoning -- four-prong no-go complete]

\noindent
\textbf{(i) Partition kills $\mathfrak{su}(3)_{\mathrm{color}}$.}
Any designation of two of the six $\mathrm{Cl}(6)$ internal directions as
spatial leaves a four-direction complement with bivector algebra
$\mathfrak{so}(4)$, $\dim 6$, which is too small to house
$\mathfrak{su}(3)$ ($\dim 8$).
[Theorem~\ref{thm:partition}, \textsc{established}]
% [Theorem 4, ESTABLISHED]

\medskip\noindent
\textbf{(ii) Double-booking triggers Coleman--Mandula.}
Not partitioning forces spatial $\mathfrak{so}(3)$ to share generators
with color $\mathfrak{su}(3)$, since the centralizer
$C_{\mathfrak{so}(6)}\!\bigl(\mathfrak{su}(3)\bigr)$ is one-dimensional
(Theorem~\ref{thm:centralizer}) and cannot house the three independent
generators $\mathfrak{so}(3)$ requires.  Coleman--Mandula's conclusion,
applied to the would-be 4D emergent theory under the standard hypotheses
(Poincar\'{e}-invariant interacting QFT with a mass gap and a nontrivial
$S$-matrix, stated in full in Section~\ref{sec:survives}), is that the
symmetry algebra is a direct sum of the Poincar\'{e} algebra and an
internal algebra commuting with it:
$[\text{Lorentz},\,\text{internal}] = 0$.  A spatial rotation generator
that shares with $\mathfrak{su}(3)_{\mathrm{color}}$ would then be
required to commute with all of Lorentz while also having nonzero
brackets with the other $\mathfrak{su}(3)$ generators
($\mathfrak{su}(3)$ is nonabelian); these cannot both hold.  The
invocation is therefore conditional: if the emergence programme reaches a
4D theory meeting Coleman--Mandula's hypotheses, that theory cannot carry
the forced spatial-color sharing, so the target is unreachable under
those hypotheses.  The escape is to void the hypotheses (Branch~A,
Section~\ref{sec:survives}).
\citep{ColemanMandula1967,HLS1975}.
% [Theorem 1, ESTABLISHED; C-M invocation conditional]

\medskip\noindent
\textbf{(iii) The chain gives only $1{+}1$.}
The external chain axis $n$ provides the $1{+}1$ Lorentzian sector and is
one-dimensional.  Three mechanisms block extension: (a)~multiple Fermi
points on a 1D chain yield parallel translations along $n$
only---different group velocities, same direction; (b)~the per-site
$M_8(\mathbb{C})$ algebra is type $\mathrm{I}_8$ (finite-dimensional,
discrete spectrum) and cannot support the HSMI machinery; (c)~the 3-leg
ladder reinterpretation of the three $\mathrm{Cl}(6)$ mode-pairs gives a
discrete transverse quantum number (3 points, no continuum limit)---the
transverse shift is a color permutation ($S_3$ Weyl element), not a
spatial translation.
\citep{MultiHSMI}.
% [ESTABLISHED]

\medskip\noindent
\textbf{(iv) The $\mathbb{H}$ factor double-books.}
The quaternion factor $\mathbb{H}$ houses $\mathrm{SU}(2)_L$ (forced by
prong~(ii); see also Section~\ref{sec:consequences}).  The $\mathbb{H}$
factor is internal by the
$\mathbb{R} \otimes \mathbb{C} \otimes \mathbb{H} \otimes \mathbb{O}$
decomposition, so its entire rotation algebra
$\mathfrak{so}(4) = \mathfrak{su}(2)_L \oplus \mathfrak{su}(2)_R$ is
internal.  The centralizer of that full internal structure is
$\operatorname{span}\{I\}$ (Theorem~\ref{thm:Hdilemma}), and the scalar
$I$ is not a rotation generator; no spatial $\mathfrak{so}(3)$ can be
sourced from~$\mathbb{H}$.
% [Theorem 5, ESTABLISHED]

\paragraph{Cross-tensor combinations.}
A candidate spatial generator need not lie wholly within one tensor
factor; it could combine a $\mathrm{Cl}(6)$ bivector part and an
$\mathbb{H}$ rotation part.  Such a generator must commute with the full
internal symmetry across both factors to avoid double-booking, so its
$\mathrm{Cl}(6)$ component must lie in
$C_{\mathfrak{so}(6)}\!\bigl(\mathfrak{su}(3)\bigr) =
\operatorname{span}\{H\}$ (prongs~(i), (ii)) and its $\mathbb{H}$
component must lie in the centralizer of
$\mathfrak{su}(2)_L \oplus \mathfrak{su}(2)_R = \operatorname{span}\{I\}$
(prong~(iv)).  The $\mathrm{Cl}(6)$ component then supplies at most one
dimension of rotation (the $\mathrm{U}(1)$ generated by $H$ in the
internal sector); the $\mathbb{H}$ component contributes no rotation
generator at all ($I$ is symmetric, not in $\mathfrak{so}(4)$).  Together
they span at most a two-dimensional abelian Lie algebra, while spatial
$\mathfrak{so}(3)$ requires a three-dimensional nonabelian one.  No
admissible combination yields $\mathfrak{so}(3)$; cross-tensor
combinations reduce to the per-factor obstructions and open no new route.

\medskip
\emph{All four prongs are computationally verified at machine precision
across three independent scripts}
(\texttt{cl6\_su3\_centralizer\_check.py},
\texttt{cl6\_spatial\_direction\_carving.py},
\texttt{cl6\_H\_spatial\_route\_check.py}).

\begin{corollary}[Verification record]
\label{cor:verification}
The following items are verified at machine precision:
\end{corollary}

\begin{table}[ht]
\centering
\small
\begin{tabular}{@{}p{5.5cm} l l@{}}
\toprule
Result & Script & Residual / value \\
\midrule
Centralizer of $\mathfrak{su}(3)$ in $\mathfrak{so}(6)$
  $= \operatorname{span}\{H\}$, $\dim 1$
  & \texttt{centralizer\_check} & $\dim = 1$ PASS \\
Alignment of centralizer with $H$
  & \texttt{centralizer\_check}
  & $|\langle v,H\rangle| = 1.0000000000$ \\
$\mathfrak{su}(3)$ octet closure
  & \texttt{centralizer\_check} & $6.20 \times 10^{-16}$ \\
Within-pair: $2/3$ adjoint $+ 1/3$ centralizer
  & \texttt{direction\_carving} & exact fractions \\
Cross-plane: $1/2$ adjoint $+ 1/2$ coset
  & \texttt{direction\_carving} & exact fractions \\
Case~A (within-pair spatial): $\mathfrak{su}(3)$ broken
  & \texttt{direction\_carving} & $1/8$ generators \\
Case~B (cross-pair spatial): $\mathfrak{su}(3)$ broken
  & \texttt{direction\_carving} & $0/8$ generators \\
All 8 cross-pair variants: $\mathfrak{su}(3)$ broken
  & \texttt{direction\_carving} & uniform $0/8$ \\
$\mathfrak{su}(2)_L$ closure in $\mathbb{H}$
  & \texttt{H\_spatial\_route} & $0.00 \times 10^{+00}$ \\
$\mathfrak{su}(2)_R$ closure in $\mathbb{H}$
  & \texttt{H\_spatial\_route} & $0.00 \times 10^{+00}$ \\
Cross-commutativity $[L_a, R_b] = 0$
  & \texttt{H\_spatial\_route} & $0.00 \times 10^{+00}$ \\
Centralizer$(\mathfrak{su}(2)_L) = \mathfrak{su}(2)_R \oplus
  \mathbb{R}$, $\dim 4$
  & \texttt{H\_spatial\_route} & $\dim = 4$ PASS \\
$\mathfrak{su}(2)_R$ nonabelian
  & \texttt{H\_spatial\_route}
  & $\max\|[R,R]\| = 4.0$ \\
Centralizer$(\mathfrak{su}(2)_L \oplus \mathfrak{su}(2)_R) =
  \operatorname{span}\{I\}$, $\dim 1$
  & \texttt{H\_spatial\_route} & $\dim = 1$ PASS \\
\bottomrule
\end{tabular}
\caption{Computational verification record.  All residuals at or below
  machine precision.  Script names abbreviated from full
  \texttt{cl6\_\ldots.py} filenames.}
\label{tab:verification}
\end{table}

%======================================================================
% SECTION 7: WHAT SURVIVES
%======================================================================
\section{What Survives}
\label{sec:survives}

\subsection{The two surviving paths}

The four-prong no-go closes all algebraic routes to $(3{+}1)$ under
Stance~1.  Two paths survive.

\paragraph{Path~A: beta (physical escape).}
The no-go prongs~(ii) and~(iv) invoke Coleman--Mandula in the standard
form, which requires: (1)~a Poincar\'{e}-invariant 4D QFT, (2)~a
nontrivial analytic $S$-matrix with a mass gap, and (3)~a finite number
of particle species.  The algebraic obstructions established here assume
all three hypotheses hold.  If the emergent infrared theory violates
hypothesis~(2)---by being genuinely gapless with no asymptotic particle
$S$-matrix---Coleman--Mandula's conclusion does not apply before its
premises do.

The committed specific escape is Branch~A: the emergent IR has no
asymptotic-particle $S$-matrix in the genuine conformal or
continuous-spectrum sense.  Confinement does not qualify---it relabels
asymptotic states as hadrons but retains an $S$-matrix; the escape
requires eliminating asymptotic states entirely.  Branch~A was committed
prior to this paper~\citep{C5gamma}; the four-prong no-go sharpens its
target by specifying exactly what the emergent IR must avoid producing: a
standard Poincar\'{e} representation on asymptotic states acting
independently of color.

The MZ and AD
theorems~\citep{MaldacenaZhiboedov2011,AlbaDiab2015} are the natural
analogues of Coleman--Mandula for CFTs: a unitary CFT in $d \geq 3$ with
a conserved higher-spin current of spin $> 2$ and a unique local stress
tensor is necessarily free.  Neither theorem engages on the C5-delta
$1{+}1$ sector, and the primary reason is dimensional.  MZ's argument is
stated for $d = 3$ and AD extends it to $d > 3$.  The $d = 2$ case
admits a richer higher-spin structure ($W_N$ algebras, integrable
theories) that the proofs of MZ and AD do not engage.  The proofs rely
on the finite dimensionality of the conformal algebra
$\mathfrak{so}(d,2)$, which fails in $d = 2$ where the conformal
symmetry is the infinite-dimensional Virasoro algebra.  The $1{+}1$
chiral sector is therefore outside the stated domain of both theorems by
dimension alone, and this is the load-bearing escape.

A second structural observation reinforces the dimensional escape but is
not required by it.  MZ and AD both take a unique local stress tensor as
an explicit hypothesis, and both sets of authors explicitly disclaim the
no-stress-tensor case.  The natural stress-energy carrier in a chiral net
of the C5-delta type is the modular Hamiltonian $K$ and its boundary form
rather than a local pointwise stress tensor in MZ's sense; whether the
emergent IR genuinely lacks a unique local stress tensor is a property of
the construction that has not been established here and is logged as
open.
% [OPEN/CONJECTURE]
The dimensional escape does not depend on it.

The genuine open question for Branch~A is whether a $d = 2$ analogue of
these results---a chiral-algebra or Virasoro-based rigidity theorem---constrains the higher-spin content of the emergent sector in a way the
$d \geq 3$ theorems do not.  Because $d = 2$ conformal symmetry is
infinitely richer than its higher-dimensional counterpart, such an
analogue would be a structurally different result, and whether one exists
that bears on the emergent sector is not known.  That reconnaissance has
not been run.  If such a result exists and bears on the emergent sector,
it would close the no-stress-tensor sub-escape before any $(3{+}1)$
extension is attempted, collapsing Branch~A onto the dimensional escape
alone; if it does not, the sub-escape stands as a second reinforcement.
Either way the dimensional escape (the load-bearing defense) is
unaffected.  Beta remains open.
% [OPEN]

A further constraint shapes what Branch~A's target could be if the
programme ever reached four dimensions.  The MZ and AD theorems engage
only a CFT carrying an exactly conserved higher-spin current of spin
greater than two; absent such a current the theorems are vacuous (the
higher-spin current is their engine, not a side hypothesis).  A 4D
emergent CFT could therefore evade them in one of two ways.  First, by
carrying no exactly conserved higher-spin current at all: but a generic
interacting CFT already satisfies this, so this route does not by itself
select an emergence target, and the relevant pressure becomes the
slightly-broken case (MZ Section~7), where higher-spin currents with
order-$1/N$ anomalous dimensions still impose constraints in the
large-$N$ regime of the $O(N)$ vector models.  Second, by lacking a
unique local stress tensor: the disclaimed case discussed above.  The one
option that is closed is a free CFT\@.  MZ forces a CFT with exactly
conserved higher-spin symmetry to free-field correlators, and a free
theory is not a viable emergence target for a programme that requires
interactions.  Branch~A's target, if it exists, is therefore an
interacting 4D theory that is either not a standard local CFT (no unique
local stress tensor) or carries only approximate higher-spin structure
subject to the slightly-broken constraints.  Which of these the emergence
programme could produce is not determined here; it is the substance of
the open beta question.

\paragraph{Path~B: Stance~2 (algebraically grounded import).}
If beta fails, the centralizer theorem provides a positive result: every
import required by Stance~2 is now algebraically justified rather than
assumed by default.

\begin{itemize}
\item Spacetime (flat GTG background): forced by prongs~(i), (ii),
  and~(iii).
\item $\mathrm{SU}(2)_L$ on the $\mathbb{H}$ factor: forced by the
  centralizer theorem ($\dim 1 < 3$).
\item $\mathrm{SU}(3)_{\mathrm{color}}$ gauge links on a separate
  $\mathbb{C}^3$: forced by the double-booking obstruction.
\end{itemize}

The G-series papers provide the working realization of the spacetime
import: Gauge Theory Gravity
(GTG)~\citep{LasenbyDoranGull1998} builds gauge-theoretic gravity in
spacetime algebra on a flat Minkowski substrate, with position-gauge field
$\bar{h}$ and rotation-gauge field $\Omega$ as the clean factorization of
the gravitational degrees of freedom.  GTG begins precisely where
Stance~2 ends---it consumes the Minkowski vector space, local
$\mathrm{Spin}(1,3)$, and smooth structure as axiomatic inputs.
Stance~2's derivation burden is exactly the substrate GTG declines to
derive.  The target is therefore known and precise: deliver the smooth
four-dimensional space with $\mathrm{Cl}(1,3)$ structure at every point
and $\mathrm{Spin}(1,3)$ acting locally, and GTG's machinery---including
the full empirical record of GR---is inherited automatically.  Stance~2
is not hypothetical; it is already partially built.

Stance~2 is the algebraically grounded fallback if beta fails.  It is not
algebraically necessary on its own---that would require beta to be closed
against---but the centralizer theorem supplies retrospective
justification for the imports Stance~2 makes, rather than leaving them as
unmotivated assumptions.  The Furey
construction~\citep{Furey2016} for $\mathrm{SU}(2)_L$ on $\mathbb{H}$,
in particular, is algebraically mandated by the centralizer theorem
rather than merely convenient.

A concrete instance of this mandate surfaced during the Paper~3
adversarial referee review: finding B-2 noted that the candidate
$T_3 = -B_{34}$ fails to commute with $\mathfrak{su}(3)_{\mathrm{color}}$.
The centralizer theorem is the explanation---$B_{34}$ carries $2/3$
adjoint content (Lemma~\ref{lem:decomposition}) and cannot be
color-neutral.  B-2 is not a gap in Paper~3; it is the centralizer
theorem asserting itself on the electroweak sector, and it forces
$\mathrm{SU}(2)_L$ off the $\mathrm{Cl}(6)$ bivectors and onto the
$\mathbb{H}$ factor exactly as Furey's construction places it.
% [ESTABLISHED reasoning -- Stance 2 as forced and justified fallback]

The two surviving paths are not symmetric.  Path~A (beta) is genuinely
open: no mechanism is in hand for producing an $S$-matrix-free emergent
IR from the $1{+}1$ chiral net, and the 2D-MZ recon is an outstanding
prerequisite before Branch~A's sub-escape can be evaluated on its merits.
Path~B (Stance~2) is already partially realized: GTG provides the
spacetime import, the Furey construction provides
$\mathrm{SU}(2)_L$ on $\mathbb{H}$, and the centralizer theorem provides
the algebraic justification for each choice.  If beta fails, the
programme does not start over---it inherits a working framework whose
imports are now explained rather than assumed.

\subsection{The hypercharge exception and its uniqueness}

An important positive consequence of Theorem~\ref{thm:centralizer} is the
uniqueness of hypercharge.  The Standard Model contains exactly one
$\mathrm{U}(1)$ factor commuting with $\mathfrak{su}(3)_{\mathrm{color}}$
(within the electroweak sector before symmetry breaking).  The centralizer
theorem shows this is not a free parameter of the model but a structural
fact of the $\mathrm{Cl}(6)$ algebra: there is room for exactly one such
$\mathrm{U}(1)$, and it is identified uniquely as $Y \propto H =
\tfrac{3}{2}\,Q_{B\text{-}L}$.

\subsection{The $1{+}1$ result is unaffected}

The partition no-go and the $\mathbb{H}$ dilemma constrain extensions of
the existing $1{+}1$ sector.  They do not affect the $1{+}1$ result
itself, which rests on the C-series chain construction and the
Tomita--Takesaki modular
structure~\citep{C4consolidation,C5delta}.  The $1{+}1$ sector---a
massless chiral fermion with a clean Lorentzian signature---stands
independently.

\subsection{Connection to virtual photon compositeness}

In the $1{+}1$ sector, $d = 2$ and the transverse polarization count
$d - 2 = 0$; gauge bosons are composite (Schwinger mechanism,
exact)~\citep{Schwinger1962}.  This section connects two regimes: the
composite-gauge-boson property of the $1{+}1$ sector the C-series has
actually built, and the activation of both obstructions in the 4D target
the programme aspires to reach.

Both obstructions trace to the same source.  Section~\ref{sec:consequences}
(the $B$--$L$ reductio) and Section~\ref{sec:partition} (the partition
no-go) each follow from the one-dimensional centralizer of
$\mathfrak{su}(3)$ in $\mathfrak{so}(6)$: the helicity route is closed
because the unique color-neutral bivector $H$ carries $B$--$L$ charge,
and the spatial route is closed because the same one-dimensional
centralizer cannot house $\mathfrak{so}(3)$.  In 4D both walls activate
together, both stemming from the centralizer dimension.  The stronger
statement that these two walls are literally the same wall---a single
obstruction wearing two faces---is established
in~\citep{TwoWalls}; the present paper establishes the weaker
and sufficient claim that both activate and share the centralizer as
their common source.

The composite gauge boson result in $1{+}1$ is the boundary statement;
the four-prong no-go is the obstruction to extending it to 4D.

%======================================================================
% SECTION 8: DISCUSSION
%======================================================================
\section{Discussion}
\label{sec:discussion}

The central result is negative: no algebraic route to $(3{+}1)$ exists
within the $\mathrm{Cl}(6)$ construction under Stance~1.  This is a
first-class result.  A clean no-go is as valuable as a positive
construction---it precisely delimits the problem that must be solved
(beta) and shows what Stance~2 is forced to import and why.

The four prongs correspond to the four candidate sources of additional
spatial directions in the construction, and close each in turn:

\begin{itemize}
\item Prong~(i) closes the partition route: carving two spatial
  directions out of the six $\mathrm{Cl}(6)$ internal directions.
\item Prong~(ii) closes the double-booking route: leaving the six
  directions unpartitioned and sharing generators between spatial
  $\mathfrak{so}(3)$ and color $\mathfrak{su}(3)$.
\item Prong~(iii) closes the chain-extension route: seeking additional
  directions from the external chain axis $n$ (parallel Fermi-point
  translations, finite-dimensional per-site algebra, and discrete ladder
  transverse quantum number all block extension beyond $1{+}1$).
\item Prong~(iv) closes the $\mathbb{H}$-factor route: sourcing spatial
  rotations from the quaternion tensor factor.
\end{itemize}

These three candidate sources---the $\mathrm{Cl}(6)$ directions, the
external chain axis, and the quaternion factor---give rise to four prongs
because the $\mathrm{Cl}(6)$ reservoir admits the partition /
no-partition dichotomy (prongs~(i) and~(ii)).  The claim is completeness
over these sources, not that each prong is individually irredundant in
some smaller logical sense.  Removing any one prong leaves its
corresponding source unaddressed, so all four are needed to close the
construction as it stands.  Whether a future extension of the
construction introduces a fifth candidate source is a separate question,
outside the present scope.

The dimension-theoretic argument at the core
($\dim\mathfrak{so}(4) = 6 < 8 = \dim\mathfrak{su}(3)$) is elementary
once stated, but its consequences were not immediate because the
algebraic mixing (Lemma~\ref{lem:decomposition}) obscures any would-be
clean separation.  The numerical verification across all case types (nine
bivector pairs in \texttt{cl6\_su3\_centralizer\_check.py}, ten distinct
spatial assignments in \texttt{cl6\_spatial\_direction\_carving.py}, seven
algebraic identities in \texttt{cl6\_H\_spatial\_route\_check.py})
confirms that the theorem is not an artifact of a particular basis
choice.

%======================================================================
% ACKNOWLEDGMENTS
%======================================================================
\section*{Acknowledgments}

Computations performed using scripts
\texttt{cl6\_su3\_centralizer\_check.py},
\texttt{cl6\_spatial\_direction\_\allowbreak carving.py}, and
\texttt{cl6\_H\_spatial\_route\_check.py}, generated and verified in the
TMFT Program research sessions (June~2026).

%======================================================================
% REFERENCES
%======================================================================
\bibliographystyle{plainnat}
\bibliography{four_prong_obstruction}

\end{document}
